\section{Discussion}\label{sec:discussion}

\subsection{What the Results Mean for Swarm Designers}
The synchronization--staleness phase transition (Thm.~\ref{thm:impossibility}) elevates a routine engineering parameter---how long to remember a teammate's telemetry---into a first-class scalability constraint.
Concretely, with $\tau = 50$\,ms slots and $T_{\mathrm{val}} = 2$\,s, our analysis caps a fully convergent TDMA-CBBA fleet at $N \leq 40$ drones; larger fleets require either faster slots, hierarchical scheduling (cluster-based TDMA), or validity horizons that scale as $T_{\mathrm{val}} = c\,N\tau$.
We emphasize that this constraint is \emph{structural}, not statistical: beyond the threshold no amount of additional runtime recovers consensus, because ghost pruning deterministically erases allocations faster than gossip can re-establish them.

Equally instructive is what does \emph{not} degrade: per-link PDR follows the analytic channel model independently of swarm logic (Sec.~\ref{sec:results-pdr}), and cohesion remains stable across radio conditions once assignments exist.
The failure mode is therefore specific and localized---information age versus scheduling period---which is precisely what makes it fixable.

\subsection{Relation to Prior Work}
Our results refine, rather than contradict, CBBA theory~\cite{choi2009}: under bounded-delay communication whose bound respects belief expiry, all classical guarantees carry over (Thms.~\ref{thm:convergence}--\ref{thm:worstcase}).
The phase transition complements asynchronous-CBBA analyses that assume infinite belief persistence; it shows those analyses do not survive the addition of standard forgetting mechanisms.
Within the demand-assigned MAC literature for tactical networks, our scheduler is closest in spirit to proximity-aware slot doubling; we treat that D-TDMA layer as an engineering extension (Sec.~\ref{sec:discussion-dtdma}) rather than as part of the core theoretical contribution.

\subsection{D-TDMA as an Engineering Extension}\label{sec:discussion-dtdma}
Our Dynamic TDMA scheduler lets an agent request a doubled transmission slot when it is near a threat or objective, prioritizing safety-critical updates.
We deliberately exclude this layer from the theoretical claims of this paper: slot doubling alters the frame structure and the effective belief-age bound, so the convergence and optimality guarantees of Section~\ref{sec:theory} are proven for the baseline fixed-slot scheduler only.
D-TDMA is implemented in the released \texttt{src/communication} package and can be enabled at runtime, but it was disabled in the benchmark matrix and its empirical characterization is left to future work.

\subsection{Limitations}
Several limitations bound the scope of our claims.
\emph{(i) Scale.}
Experiments span $n \leq 50$ in simulation; physical swarms introduce latency jitter, asymmetric links, and localization error absent here.
\emph{(ii) Static tasks.}
Objectives are stationary; dynamic task arrival would exercise bundle re-planning paths not covered by our single-claim auction.
\emph{(iii) Utility measurement.}
Theorem~\ref{thm:worstcase} is validated structurally (conflict-freeness at consensus); we do not yet report realized utility against a Hungarian-optimal upper bound, so Eq.~\eqref{eq:optimality-bound}'s tightness in practice remains open.
\emph{(iv) Homogeneity and attrition.}
All agents share parameters, and battery attrition is disabled in the benchmark corpus, so the fault-tolerance path of ghost pruning is exercised only through belief expiry rather than true agent loss.
\emph{(v) D-TDMA coverage.}
The D-TDMA extension is specified and implemented but was disabled in the released matrix; its empirical characterization is future work.
\emph{(vi) Loss asymmetry.}
The $p_{\mathrm{rx}} = p_{\mathrm{tx}}/2$ convention is a modeling choice; sensitivity to this ratio is unexplored.

\subsection{Future Work}
Four directions follow directly.
First, \emph{adaptive horizons}: let each agent set $T_{\mathrm{val},i} = \hat{c}\,\widehat{N}\tau$ from its estimated neighbor count, converting Theorem~\ref{thm:impossibility}'s constraint into a self-scaling protocol.
Second, \emph{utility-grounded evaluation}: add an optimal-assignment oracle to measure realized optimality ratios and empirically test Theorem~\ref{thm:worstcase}'s tightness.
Third, \emph{hierarchical scheduling}: cluster-based TDMA can shrink effective frame length below $N\tau$, pushing the phase transition outward without faster radios.
Fourth, \emph{learned slot allocation}: replace proximity-triggered doubling with a lightweight learned policy that anticipates contention, trading a small compute budget for further latency reduction.
Finally, physical-hardware trials on a $10$--$20$ drone testbed with real 2.4\,GHz radios would anchor the simulation-to-reality gap.
